\section{Results}
\label{sec:results}

Every number below comes from a stored run under \texttt{results/} and was read
back from the file named beside it. The synthetic grid holds $22$ scenarios
drawn from the three generating families, each fitted on three replicates, with
the replicate seed fixed by the scenario so that the same scenario and
replicate index reproduce the same dataset. Across the five grid experiments
this is $1176$ scored fits; the parameter ladder adds $81$ and the external
benchmark adds $120$.

One property of the generator governs how the structural tables read. Over the
$66$ (scenario, replicate) pairs of the full grid the generator places $731$
lagged edges, $482$ of them at scale $k=2$ and $249$ at scale $k=3$, so the
mean number of true edges per fit is $11.076$ and two of the three scales carry
a recovery test. Scale $1$ carries no edge in any replicate, so $\gP_{1,\ell}$
is empty for every label $\ell$. That emptiness is a measured property of the
generator rather than a convenience. The finest scale is a binary event-start
process, so a lagged edge there can only shift which minute of a block begins
an episode, which is worth less than the price of the edge, and the screen of
\cref{sec:benchmark_identifiability} rejects it in every family. A scale with
no true edge and no proposed edge is scored as perfectly recovered, with
precision, recall and $F_1$ all equal to $1$; a scale with no true edge and at
least one proposed edge scores $0$ on all three. Reported precision at scale
$1$ is therefore the fraction of fits that proposed nothing there, not a ratio
of counts, and any mean of $F_1$ across the three scales still carries that
free term. One of the $66$ pairs happens to carry no true edge at scale $3$
either, which is why the empty-graph row scores $0.015$ there rather than
exactly $0$.

\subsection{Structure recovery on the synthetic grid}
\label{sec:results:structure}

Table~\ref{tab:results:exp1} reports the ideal-recovery experiment
(\texttt{results/exp1/exp1\_results.csv}, $198$ fits), where the generator, the
number of regimes and the number of factors are all given to the learner and
only $\gP_{k,\ell}$ has to be found. Every fit on the grid
starts at the generating parameters, and the structure search starts at the
generating parent sets, so the structure reported for the full model is where
the penalized score moves from the truth rather than where a cold search
arrives; the empty-graph row clears those parent sets and runs no search.

\begin{table}[t]
\centering
\caption{Structure recovery on the synthetic grid, means over $22$ scenarios
and three replicates. Scale $1$ carries no true edge, so its precision, recall
and $F_1$ record the fraction of fits that proposed nothing there. In the
across-scales row the $F_1$ column is the mean over the three scales and the
SHD column is their sum. Source
\texttt{results/exp1/exp1\_results.csv}.}
\label{tab:results:exp1}
\begin{tabular}{llrrrr}
\toprule
Method & Scale & Precision & Recall & $F_1$ & SHD \\
\midrule
n$^2$DBN      & $k=1$ & 0.273 & 0.273 & 0.273 & 2.318 \\
              & $k=2$ & \textbf{0.864} & \textbf{0.598} & \textbf{0.680} & \textbf{3.379} \\
              & $k=3$ & \textbf{0.985} & \textbf{0.870} & \textbf{0.908} & \textbf{0.515} \\
              & across &      &       & 0.620 & 6.212 \\
\midrule
Empty graph   & $k=1$ & 1.000 & 1.000 & 1.000 & 0.000 \\
              & $k=2$ & 0.000 & 0.000 & 0.000 & 7.303 \\
              & $k=3$ & 0.015 & 0.015 & 0.015 & 3.773 \\
              & across &      &       & 0.338 & 11.076 \\
\midrule
Oracle parameters & all & 1.000 & 1.000 & 1.000 & 0.000 \\
\bottomrule
\end{tabular}
\end{table}

The total structural Hamming distance favors the full model. It averages
$6.212$ against $11.076$ for the empty graph, so proposing a structure saves
$4.864$ edge errors per fit on average against proposing nothing. Paired within
(scenario, replicate), the full model is better on $61$ pairs, tied on $5$ and
worse on none of $66$.

The decomposition locates the gain, and it falls on two scales rather than
one. At scale $2$ the paired difference is $-3.924$ edge errors per fit, better
on $58$ pairs, tied on $6$ and worse on $2$. At scale $3$ it is $-3.258$,
better on $64$ pairs, tied on $2$ and worse on none. At scale $1$, where the
truth is empty, the full model proposes on average $2.318$ edges that cannot be
correct, losing every pair in which it proposes anything at all ($0$ wins, $18$
ties, $48$ losses); it proposes nothing there on $18$ of the $66$ fits and
$3.188$ edges on each of the rest. Summed, a gain of about seven edge errors
across the two scales that carry the truth is offset by about two false
positives at the scale that carries none, so the total keeps most of the gain,
and reporting it alone would still hide all three effects.

The silence at the finest scale is selective rather than a general reluctance
to propose. The same fits propose at least one edge at scale $2$ on $63$ of
$66$ and at scale $3$ on $65$ of $66$.

Precision at scale $2$ runs $0.851$ on family I, $0.775$ on family II and
$0.959$ on family III, against recall $0.655$, $0.502$ and $0.668$. At scale
$3$ precision runs $0.917$, $1.000$ and $1.000$ against recall $0.833$, $0.836$
and $0.921$. The selection is conservative in the same way in all three
families, and the coarsest scale is the easier of the two scales that carry
true edges, returning more of the truth at a higher precision than the middle
scale in every family.

\subsection{Prediction against the empty graph}
\label{sec:results:prediction}

Held-out fit is scored by the validation loss $-\sum_k \log Z_k$ evaluated on
the validation split with the learned $\Theta$ and the learned graph, where
$\log Z_k$ is the normalizer the scale-$k$ E-step returns, the forward-backward
partition function of the regime chain evaluated with the factor posterior
folded into the emission terms. Lower is better. Because the comparison is
between two methods fitted on the same split, the paired difference is the
quantity of interest and the level is not.

Over the $66$ pairs of \texttt{results/exp1/exp1\_results.csv} the full model
reaches a mean of $-36296.45$ against $-34951.83$ for the empty graph. The mean
paired difference is $-1344.63$ nats and the median is $-632.20$; the full model
is better on $64$ of $66$ pairs. Expressed relative to the
empty-graph value of each pair, the mean reduction is $4.0$ percent. The
selected edges therefore pay for themselves in held-out evidence on the same
fits where they also reduce the total edge distance.

The oracle-parameter row records a validation loss of $-3485.14$, far above
both. That row installs the generating parameters and fits nothing, so it is
the ceiling for structure and for latent state, not for prediction; a
parameter set tuned on the training split will fit a held-out split of the same
series better than the parameters that generated it.

\subsection{Latent-state recovery against the oracle ceiling}
\label{sec:results:latent}

Table~\ref{tab:results:latent} scores the recovered regime path $S_t^{(k)}$
against the generating path, with the oracle-parameter row as the ceiling.

\begin{table}[t]
\centering
\caption{Latent regime recovery, means over the full grid. The oracle row
installs the generating $\Theta$ and runs inference only, so it bounds what any
estimator on this grid can reach. Source
\texttt{results/exp1/exp1\_results.csv}.}
\label{tab:results:latent}
\begin{tabular}{lrrrrrr}
\toprule
& \multicolumn{3}{c}{Adjusted Rand index} & \multicolumn{3}{c}{Best-permutation accuracy} \\
\cmidrule(lr){2-4}\cmidrule(lr){5-7}
Method & $k=1$ & $k=2$ & $k=3$ & $k=1$ & $k=2$ & $k=3$ \\
\midrule
n$^2$DBN          & --- & 0.295 & 0.907 & --- & 0.767 & 0.961 \\
Empty graph       & --- & 0.237 & 0.908 & --- & 0.736 & 0.967 \\
Oracle parameters & --- & 0.219 & 0.967 & --- & 0.742 & 0.987 \\
\bottomrule
\end{tabular}
\end{table}

Most of this table is empty on purpose, and that is the result rather than a
gap in it. A scale whose generating path never switches has no regime to
recover, and the screen of \cref{sec:benchmark_identifiability} places a regime
only where its own oracle clears the constant predictor, so on the minute grid
scales $1$ and $2$ carry one state and are reported as not applicable rather
than as a perfect score. The scale-$2$ column holds the three day-grid fits,
whose hour scale does carry a regime, and the scale-$3$ column holds $25$ of the
$27$ family II fits, the two omissions being replicates whose coarsest test
partition happened not to switch.

Where a regime is placed, it is recovered. At the coarsest scale the full model
reaches an adjusted Rand index of $0.907$ and an accuracy of $0.961$ against an
oracle at $0.967$ and $0.987$, and it recovers the path exactly on $17$ of its
$25$ scored fits, the same count as the oracle. The gap to the ceiling is
$0.060$ of index and $0.026$ of accuracy. At the hour scale of the day grid the
three rows are $0.295$, $0.237$ and $0.219$ of index, where the fitted model is
ahead of the oracle on three fits, which is inside the noise of three fits and
should not be read as beating a ceiling.

An earlier version of this benchmark placed a regime at all three scales and
reported near-identical rows for the full model, the empty graph and the oracle,
with scale-$2$ accuracies of $0.529$, $0.521$ and $0.516$ against a floor of one
half. That was not three methods tying; it was three methods scoring a quantity
the data did not identify. The screen removed those scales from the grid, and
what is left is the one comparison that carries information.

The latent factor $H_t^{(k)}$ tells the same story with a smaller gap. On the
family III scenarios the mean absolute correlation between the recovered and
generating factor trajectories is $0.457$, $0.470$ and $0.694$ at the three
scales for the full model, against $0.540$, $0.555$ and $0.786$ for the oracle.
The estimator gives up between $0.083$ and $0.092$ of correlation at each
scale, and the ceiling itself is well below one.

\subsection{The contribution of the cross-scale coupling}
\label{sec:results:coupling}

The coupling experiment (\texttt{results/exp2/exp2\_results.csv}, $132$ fits)
removes the downward pass, so that the regime chain of scale $k$ is driven only
by that scale's own data and context $\gK_t^{(k)}$ and never by the posterior
of its parent epoch $\parmap_k(t)$. Everything else, including the allocation
that ties $X_t^{(k)}$ and $Y_t^{(k)}$ to their parent totals, is unchanged.

On the minute grid the coupling now buys exactly nothing, and the reason is
structural rather than empirical. Its sixty fits return identical graphs and
identical validation losses to the last recorded digit, on every column. With
$M_k = [1,1,M]$ the parent posterior passed downward at scales $1$ and $2$ is a
distribution over one state, which is the same object whether or not it is
passed, so the coupled and uncoupled arms are the same computation. The
experiment does not test the coupling there.

The six day-grid fits are the ones that test it, since the hour scale of that
hierarchy does carry a regime with a two-state parent above it. There the
uncoupled arm is ahead on held-out evidence by a mean of $362.5$ nats, with the
coupled model better on $2$ of the $6$ pairs; the structural columns are
identical. Six fits with that spread support no claim in either direction, so
the honest statement is that this benchmark no longer measures the contribution
of the cross-scale coupling.

This is a direct cost of the screen, and it is worth stating as one. An earlier
version of this grid put a regime at all three scales and recorded the coupling
as worth $61.89$ nats and $0.061$ edge errors, all of it on family II. What the
coupling was passing downward there was a finer-scale regime whose own oracle
could not beat a constant predictor, so that figure measured the benchmark
rather than the mechanism. Removing the unrecoverable latent removed the
measured benefit with it. Testing the coupling properly needs a hierarchy in
which two adjacent scales both carry a latent the screen admits, which on the
present grid is one hierarchy out of two and three scenarios out of twenty-two.

\subsection{Component ablation}
\label{sec:results:ablation}

The ablation (\texttt{results/exp3/exp3\_results.csv}, $270$ fits over the $18$
family II and family III scenarios) removes one component at a time. Dropping
the factor sets $r_k = 0$; dropping the regime sets $M_k = 1$.

\begin{table}[t]
\centering
\caption{Component ablation on the family II and family III scenarios, means
over $18$ scenarios and three replicates. Source
\texttt{results/exp3/exp3\_results.csv}.}
\label{tab:results:ablation}
\small
\begin{tabular}{lrrrrrrr}
\toprule
Method & $F_1$ $k{=}2$ & $F_1$ $k{=}3$ & SHD $k{=}1$ & SHD $k{=}3$ & SHD all & Acc.\ $k{=}3$ & Val.\ loss \\
\midrule
n$^2$DBN          & \textbf{0.677} & \textbf{0.918} & 2.362 & \textbf{0.500} & \textbf{6.259} & \textbf{0.961} & \textbf{-37814.78} \\
No latent factor  & 0.677 & 0.918 & 2.352 & 0.537 & 6.259 & 0.961 & -37767.39 \\
No regime chain   & 0.676 & 0.784 & 2.389 & 1.796 & 7.574 & 0.376 & -37692.81 \\
Empty graph       & 0.000 & 0.000 & 0.000 & 3.815 & 11.130 & --- & -36454.00 \\
Oracle parameters & 1.000 & 1.000 & 0.000 & 0.000 & 0.000 & 0.987 & -4033.88 \\
\bottomrule
\end{tabular}
\end{table}

The two components contribute unequally, and one of them contributes nothing
structural. Removing the factor leaves the total distance equal to three
decimals, equal on $45$ of the $54$ pairs and better on $4$ against worse on
$5$, and leaves scale-$2$ and scale-$3$ $F_1$ unchanged to three decimals. Its
whole contribution is $47.39$ nats of held-out evidence, identical on $27$
pairs, better on $24$ and worse on $3$. Removing the regime chain costs
$121.97$ nats, better on all $27$ pairs that differ and worse on none, and it
also costs structure at the coarsest scale. The scale-$3$ distance rises from
$0.500$ to $1.796$, scale-$3$ $F_1$ falls from $0.918$ to $0.784$, the total
distance rises from $6.259$ to $7.574$ and the mean $F_1$ across scales falls
from $0.606$ to $0.561$.

Scale-$2$ $F_1$ alone does not show this. It is $0.676$ without the regime
chain against $0.677$ with it, so the regime pays for itself at the coarsest
scale and not at the middle one. That is where persistent drift was always
claimed to live, and it is now the only scale at which the screen of
\cref{sec:benchmark_identifiability} admits a regime on this grid, which is why
the regime's held-out contribution is $122$ nats here where an earlier version
of this benchmark recorded $1250$. That earlier figure credited the model for
tracking a finer-scale regime whose own oracle could not beat a constant, so it
measured the benchmark rather than the method. The regime accuracy column makes
the same point from the other side: with the chain removed the fitted path is
constant, and the $0.376$ it scores is the frequency of the modal true regime,
which is the floor rather than a recovery. The regime chain is what makes the
coarsest scale recoverable, and the factor is a small likelihood term that the
structure search does not see. That is an honest negative for the factor as a
structural device.

\subsection{Behavior under misspecification}
\label{sec:results:misspec}

The misspecification experiment
(\texttt{results/exp4/exp4\_results.csv}, $36$ fits) generates from a family
that violates a modeling assumption and fits the same model. The
\emph{heavier tail} arm draws counts from a distribution with a heavier tail
than the fitted emission; the \emph{omitted context} arm generates with a
context that is then withheld from $\gK_t^{(k)}$ at fit time. No fit failed and
no run raised an error.

\begin{table}[t]
\centering
\caption{Behavior under misspecification, means over four scenario shapes and
three replicates per arm. Structure metrics only, since the three arms generate
from different data and their held-out losses are not on a common scale. Source
\texttt{results/exp4/exp4\_results.csv}.}
\label{tab:results:misspec}
\small
\begin{tabular}{lrrrrrrr}
\toprule
Arm & Prec.\ $k{=}2$ & Rec.\ $k{=}2$ & $F_1$ $k{=}2$ & $F_1$ $k{=}3$ & SHD $k{=}1$ & SHD $k{=}2$ & Sweeps \\
\midrule
Well specified   & 0.965 & 0.578 & 0.707 & 1.000 & 3.000 & 3.333 & 6.33 \\
Heavier tail     & 0.875 & 0.611 & 0.710 & 0.976 & 3.583 & 3.917 & 6.75 \\
Omitted context  & 0.846 & 0.536 & 0.650 & 0.889 & 3.667 & 4.500 & 6.17 \\
\bottomrule
\end{tabular}
\end{table}

The two arms fail in different directions, and they reach different scales. A
heavier tail than the emission assumes trades precision for recall at the middle
scale, precision falling from $0.965$ to $0.875$ while recall rises from $0.578$
to $0.611$, so its $F_1$ there is level at $0.710$ against $0.707$; what it costs
is the coarsest scale, whose $F_1$ falls from an exact $1.000$ to $0.976$, and
the empty finest scale, where false positives rise from $3.000$ to $3.583$.
Withholding the context costs both halves of the middle scale, precision
$0.846$ and recall $0.536$, and it is the arm that reaches the coarsest scale
hardest, where precision falls to $0.917$ and recall to $0.875$. Neither arm
breaks the model. Both keep the aggregation exact, at a reconciliation error of
exactly $0$ for counts and durations on all $36$ fits, and neither takes
materially longer to converge than the well-specified arm, at $6.75$ and $6.17$
outer sweeps against $6.33$.

The mean of $F_1$ across scales is the wrong summary for this table and is left
out of it. That mean is $0.590$ under the heavier tail against $0.569$ under the
well-specified arm, reversing the ordering of the coarsest scale, and the
reversal comes from the empty finest scale, where the heavier-tail arm proposes
nothing on one of its twelve fits and the other two arms on none of theirs. The
heavier tail is not the easier setting.

\subsection{Against eight published structure learners}
\label{sec:results:comparison}

The comparison experiment (\texttt{results/exp5/exp5\_results.csv}, $540$ fits)
runs eight published structure learners on the same counts $X_t^{(k)}$ that
the local score sees, then fits this model's parameters on whatever graph each
returns, so that every row is scored by the same metric code on the same data.

\begin{table}[t]
\centering
\caption{Eight published structure learners against the latent-conditional
local score and against reporting nothing, means over $18$ scenarios and three
replicates. Sorted by total structural Hamming distance. Source
\texttt{results/exp5/exp5\_results.csv}.}
\label{tab:results:comparison}
\small
\begin{tabular}{lrrrrrrr}
\toprule
Method & Prec.\ $k{=}2$ & Rec.\ $k{=}2$ & $F_1$ $k{=}2$ & $F_1$ $k{=}3$ & SHD all & $\times$ empty & Val.\ loss \\
\midrule
n$^2$DBN local score & \textbf{0.867} & \textbf{0.585} & \textbf{0.677} & \textbf{0.918} & \textbf{6.259} & \textbf{0.56} & \textbf{-37814.78} \\
Empty graph          & 0.000 & 0.000 & 0.000 & 0.000 & 11.130 & 1.00 & -36454.00 \\
Chow-Liu tree        & 0.210 & 0.196 & 0.195 & 0.082 & 23.296 & 2.09 & -36933.15 \\
PC                   & 0.073 & 0.129 & 0.088 & 0.031 & 26.815 & 2.41 & -36784.06 \\
GES                  & 0.146 & 0.220 & 0.169 & 0.043 & 27.685 & 2.49 & -36960.92 \\
Granger, BIC         & 0.164 & 0.351 & 0.217 & 0.341 & 29.389 & 2.64 & -37517.01 \\
Hill climbing, BIC   & 0.139 & 0.288 & 0.183 & 0.043 & 30.574 & 2.75 & -37452.14 \\
Exhaustive, BIC      & 0.139 & 0.288 & 0.183 & 0.043 & 30.574 & 2.75 & -37452.19 \\
DYNOTEARS            & 0.155 & 0.327 & 0.203 & 0.363 & 30.759 & 2.76 & -37556.48 \\
Hill climbing, BDeu  & 0.136 & 0.288 & 0.181 & 0.096 & 33.241 & 2.99 & -37454.06 \\
\bottomrule
\end{tabular}
\end{table}

All eight land between $2.1$ and $3.0$ times worse than reporting nothing.
Their $F_1$ at the middle scale runs from $0.088$ to $0.217$ against $0.677$
for the local score, and no fit of any of the eight beats the empty graph on
total distance in any of the $54$ pairs. The local score is the only entry
below the empty graph, by a mean of $4.870$ edge errors with $50$ wins, $4$ ties
and no losses.

The coarsest scale is where the separation is widest, and it is a column this
comparison could not carry until the grid placed a true edge there. Six of the
eight score below $0.100$ at scale $3$, and the two that do not, DYNOTEARS at
$0.363$ and the Granger test at $0.341$, still reach less than four tenths of
the local score's $0.918$. Those two are the only entries with an explicitly
temporal search. The other six run from $0.031$ to $0.096$ at that scale, near
enough to the empty graph's $0.000$ that what they mostly contribute at the
hourly scale is false positives.

\paragraph{The mechanism is a lagged self edge.}
The generator drives every label of a scenario from a shared regime $S_t^{(k)}$
that persists across epochs, so every label it drives is autocorrelated whether
or not any lagged edge points into it. A selection criterion that conditions on
nothing sees that autocorrelation and has one place to put it, namely an edge
from the label to itself at lag one. Splitting recall at scale $2$ by whether
the true edge points from a label to itself or from one label to another states
this directly rather than by inference. The Granger test recovers $1.000$ of
the true self edges and $0.210$ of the true cross edges; DYNOTEARS recovers
$0.986$ against $0.190$, and both hill climbing variants and exhaustive search
$0.957$ against $0.167$; PC recovers $0.729$ against $0.002$ and GES $0.543$
against $0.165$. The Chow-Liu tree is the one exception and is weak on both,
$0.429$ and $0.167$. The local score recovers $0.929$ of the self edges and
$0.524$ of the cross edges, and is the only entry of the nine above a quarter on
the cross edges. At the coarsest scale it is $0.944$ and $0.850$, so once the
regime is conditioned on the two kinds of edge are recovered at nearly the same
rate.

The proposed edge sets say the same thing from the other side. Across the $54$
(scenario, replicate) pairs of this experiment the generator places $391$ true
edges at scale $2$, of which $61$ point from a label to itself, a true
self-edge share of $0.156$. The proposed share is $0.985$ for PC, $0.886$ for
DYNOTEARS, $0.883$ for hill climbing with BIC and for exhaustive search,
$0.879$ for hill climbing with BDeu, $0.838$ for the Granger test, $0.743$ for
GES and $0.634$ for the Chow-Liu tree, against $0.280$ for the local score. Six
of the eight place more than four fifths of their scale-$2$ budget on a
construct that is about a sixth of the truth. Conditioning the score
on $q(\gS,\gH)$ removes that autocorrelation before an edge is priced, which is
what the precision column of Table~\ref{tab:results:comparison} records.

\paragraph{Held-out likelihood does not order the structures.}
Every one of the nine non-empty rows of
Table~\ref{tab:results:comparison} improves on the empty graph in validation
loss, by between $330.1$ and $1360.8$ nats, and the ordering it induces is close
to unrelated to the structural one. Spearman rank correlation between
mean total distance and mean validation loss across the ten rows is $-0.333$.
The empty graph has the worst validation loss of the ten and
the second-best distance; DYNOTEARS has the second-best validation loss and the
ninth-best distance. The local score leads both columns here, which is a
property of that row rather than of the criterion, since the other nine rows
are ordered almost independently by the two. Held-out likelihood is a real
quantity and it is not a proxy for structure recovery on this grid, so a paper
that reports only one of the two can support a conclusion the other refutes.

\subsection{Parameter recovery on a sample-size ladder}
\label{sec:results:params}

Structure recovery says which edges came back and nothing about the conditional
distributions attached to them. The parameter experiment
(\texttt{results/exp7/exp7\_results.csv}, $81$ rows) scores the learned
$\Theta$ against the generating $\Theta$ directly, in identified form, so that
two parameter sets inducing the same distribution score zero. It holds the
hierarchy shape fixed and varies only $T_1$, the number of finest-scale epochs,
over $10{,}080$, $20{,}160$ and $40{,}320$, chosen so that the coarsest scale
receives $168$, $336$ and $672$ epochs and the structural benchmark's own
horizon sits in the middle. Two starting points are separated. The
\emph{oracle start} begins the estimator at the generating parameters and
measures the estimator alone; the \emph{data start} builds its initial
parameters from moments of the data and measures the estimator together with
its initializer, which is the only setting that transfers to a log with no
ground truth.

The criterion was set before the study ran. Parameter recovery is a consistency
claim, so the error has to fall as $T_1$ grows; an error that does not fall
with sample size is a bias rather than an estimate.

\begin{table}[t]
\centering
\caption{Mean identified parameter error against $T_1$, over three generating
families and three replicates. The ladder is set from the coarsest scale rather
than the finest, so these three rungs give it $168$, $336$ and $672$ epochs;
the middle rung is the horizon the structural benchmark itself runs at. The third stored method installs the generating
parameters and fits nothing, so its error is zero at every rung and it is not
shown. The error averages whichever scale-free blocks a family identifies,
which is $5$ for family I, $6$ for family II and $8$ for family III, so a level
is comparable along a row and not down the family rows; only the per-family rows
carry a criterion. Source \texttt{results/exp7/exp7\_results.csv}.}
\label{tab:results:ladder}
\begin{tabular}{lrrr}
\toprule
Start & $T_1 = 10080$ & $T_1 = 20160$ & $T_1 = 40320$ \\
\midrule
Oracle start (estimator alone) & 0.0725 & 0.0695 & \textbf{0.0514} \\
Data start (estimator and initializer) & 0.2822 & 0.2595 & \textbf{0.2504} \\
\midrule
Family I, data start   & 0.3366 & 0.2735 & 0.3001 \\
Family II, data start  & 0.2145 & 0.2119 & \textbf{0.1924} \\
Family III, data start & 0.2956 & 0.2932 & \textbf{0.2588} \\
\bottomrule
\end{tabular}
\end{table}

Both starts meet the criterion, and the reading of this table changed when the
benchmark did. From the oracle start the mean identified parameter error goes
$0.0725$, $0.0695$ and $0.0514$, falling at every rung. From the data start it
goes $0.2822$, $0.2595$ and $0.2504$, also falling at every rung. Per family the
data start falls at every rung for family II, $0.2145$, $0.2119$ and $0.1924$,
and for family III, $0.2956$, $0.2932$ and $0.2588$. Family I is the one that
does not, $0.3366$, $0.2735$ and $0.3001$, and it carries neither a regime nor a
factor, so its reported error is the coarsest-scale emission block alone.

An earlier version of this grid placed a regime at all three scales and reported
a data start rising at every rung, $0.4261$, $0.4533$ and $0.4678$, with family
II alone at $0.6461$, $0.7933$ and $0.8445$. That rise is what the screen of
\cref{sec:benchmark_identifiability} removed. Three of family II's scored blocks
were the finest-scale transition table and the two allocation shares, all driven
by a regime whose own oracle could not beat a constant predictor, so the quantity
that grew with the sample size was the noise in an unidentified block rather than
a fit converging elsewhere. With the regime placed only where it is recoverable,
family II falls to about a quarter of its former level and becomes the best of
the three families rather than twice the worst. Families I and III are unchanged
to the digit, because a scenario carrying one regime at every scale is untouched
by where a regime is placed, which is the control this comparison has that a
rerun on a freshly randomized grid would not.

What remains is the gap between the two starts rather than a failure of either.
At the longest rung the data start sits at $0.2504$ against $0.0514$ from the
oracle start, so the initializer, not the estimator, is what the remaining error
measures. The coarsest transition table is recovered closely from either start,
at $0.0242$, $0.0417$ and $0.0167$ in total variation from the data start, and
the coarsest lagged support is recovered at $4.00$, $4.22$ and $5.11$ selected
keys against a true $4.00$, so the over-selection that remains is at the longest
rung and is small.

\paragraph{The objective and the parameter error no longer disagree.}
Comparing the two starts on their final objective rather than on their parameter
error separates a fit that found a worse optimum from one that found an equally
good optimum at different parameters, and the two have different remedies. Over
the $27$ paired fits the data start reaches a strictly worse objective in $3$ of
the $9$ family I fits, $2$ of $9$ for family II and none for family III, and in
every one of those it is worse on the parameter error as well. There is no fit
in which a worse objective comes with better parameters, where an earlier version
of this grid had $4$ of $27$. That pathology was the sign that the error was
being scored partly on quantities the likelihood does not see, and it is gone
with them. Two conditions attach: equality is equality at a relative tolerance,
not proven equality, and the objective compared is the monitored one, which
Section~\ref{sec:lim_meanfield} records as contained rather than exact.

\paragraph{The coarsest scale is recoverable from a data start, and only with
the level correction.}
Counting the lagged coefficient keys each fit ends with separates a coefficient
that is wrong from one that was never estimated. Against a true $4.00$ keys per
fit at the coarsest scale, a data start recovers $4.00$, $4.22$ and $5.11$
across the three rungs, so the support is recovered and the over-selection that
remains is at the longest rung. Removing two things restores the earlier behavior
exactly: initializing the coarsest level from the observed mean, which already
contains the contribution of the lagged parents, and letting the parameter step
refit that level against the empty graph before the structure search runs. With
both removed the same fits recover $0.00$, $0.00$ and $0.67$. The level is
therefore not a nuisance parameter here. Fitted before the search and against a
graph with no edges in it, it absorbs exactly what a true lagged edge would
explain, and the search then meets a model that already accounts for the data.

\paragraph{The short coarse series was the likeliest explanation, and it is
ruled out.}
The ladder holds the hierarchy shape fixed and varies $T_1$ alone, so the
coarsest scale receives $168$, $336$ and $672$ epochs across the three rungs, of
which a fit sees $117$, $235$ and $470$ after the training split, while the
finest receives $10{,}080$, $20{,}160$ and $40{,}320$. That scale now carries lagged
coefficients of its own, which made it the natural suspect. Reading the error
per scale and per parameter block removes it from suspicion on three counts.
The lagged coefficients of any scale are excluded from the reported error by
construction, since it averages the scale-free blocks alone, and perturbing the
coarsest lagged coefficients of a set of true parameters leaves it unchanged to
machine precision. Every coarsest-scale block falls at every rung from the
oracle start, the relative emission error $0.603$, $0.562$, $0.350$ and the
transition distance $0.224$, $0.074$, $0.068$. And recomputing the reported
error over subsets of the scales separates the readings directly: the coarsest
scale alone falls, $0.490$, $0.448$, $0.271$, while dropping it leaves the
ladder broken, $0.416$, $0.529$, $0.248$. A ladder that grew every scale at the
same rate would therefore not repair this one.

What the breakdown puts in its place is narrower and less comfortable. From the
oracle start the non-monotonicity is confined to family II, whose level is set
by the relative emission error at the \emph{finest} scale, reaching $7.9$ and
$9.5$ on its worst fits against $0.2$ at the coarsest, while families I and III
fall at every rung. That is the binary event-start scale on which the leak of
\cref{sec:lim_leak} sits and which the screen of
\cref{sec:benchmark_identifiability} places below its identifiability floor. The
residual therefore concentrates where this benchmark cannot separate a fitting
failure from a scale it was never able to identify, so the criterion stands as
failed and the cause is attributed no further than that.

\paragraph{The leak also shows up in the coefficients.}
The same rows record the absolute learned lag-coefficient mass placed on lagged
keys outside the true support, which is the parameter-space counterpart of a
false-positive edge. That mass is largest at the longest rung from either start
instead of shrinking. From the oracle start it grows at every step, from
$0.281$ at $T_1 = 1440$ to $1.194$ at $T_1 = 2880$ and $1.585$ at
$T_1 = 5760$; from the data start it runs $0.278$, $0.860$ and $1.376$ over the
same three rungs. The share of fits carrying any off-support mass runs $0.111$,
$0.444$, $0.556$ from the oracle start and $0.111$, $0.333$, $0.444$ from the
data start. It is not confined to one family. At the longest rung from the
oracle start the mass per fit is $2.384$ on family I, $1.458$ on family II and
$0.913$ on family III. It is confined to one scale. Of the $1.585$ overall at
that rung, $1.358$ sits at scale $1$, which is the scale the generator leaves
empty. The co-movement a persistent regime induces between
labels is being absorbed into
lagged coefficients, and a longer series makes that absorption more confident,
not less. The structural score suppresses it as an edge; the parameter step
still pays for it as a coefficient.

\subsection{External ground truth}
\label{sec:results:external}

Every number to this point is scored against data this benchmark generated,
which is the weakest evidence base for a claim about a learner. The external
experiment (\texttt{results/exp8/exp8\_structure.csv} and
\texttt{results/exp8/exp8\_cpt.csv}) repeats the structure comparison against
released networks. The bnlearn repository \citep{scutari2010bnlearn}
distributes fully parameterized networks with an exact directed graph and exact
conditional probability tables, and pgmpy \citep{ankan2024pgmpy} vendors them
inside its wheel, so the experiment runs with no network access. A static
network becomes a two-slice temporal network by mapping each edge onto a lag-1
edge, the standard construction with published precedent
\citep{trabelsi2013benchmarking, ouyang2023divideconquer}. Five networks (asia,
cancer, child, earthquake and sachs, from $5$ to $20$ nodes and $4$ to $25$
edges), three replicates and $4000$ epochs give $120$ structure fits including
the empty-graph reference, with no errors. The largest network of the earlier
version of this experiment, insurance at $27$ nodes, is not included, because
selecting the free quantities per dataset multiplies the number of fits by
eleven and that run did not complete. The omission is of size rather than of an
inconvenient result.

A tiled network is flat, so a two-scale hierarchy is imposed by summing blocks
of consecutive epochs and only the finest scale is scored, where the candidate
set is exactly the tiled ground truth. The block length is declared as the
square root of the series length, for the reason given in
\cref{sec:results:selection}, and the regime count, the factor dimension and
the edge penalty are selected per dataset on held-out log evidence.

Structure and parameters are scored on the same graph. A learned parent set
that differs from the true one produces a table of a different shape, so tables
cannot be compared entry by entry across methods. Both the learned model and
the released network do answer one common question, namely the distribution of
a node at epoch $t$ given the whole previous epoch, so that conditional is what
is scored: the table is estimated from the sampled data on the learned parents
alone, the released tables are read at the true parent configuration, and the
two are compared in total variation, averaged over epochs. An empty parent set
is a valid answer under this score and is charged the distance from the
marginal to the truth.

\begin{table}[t]
\centering
\caption{Structure and conditional probability recovery against released
bnlearn networks tiled into two-slice temporal networks, means over five
networks and three replicates. Same structural metric code as
Table~\ref{tab:results:comparison}, external ground truth. The last column is
the mean total variation distance of the conditional each method's own learned
parent set implies, so structure and parameters are scored on one graph.
Source \texttt{results/exp8/exp8\_structure.csv}.}
\label{tab:results:external}
\begin{tabular}{lrrrrr}
\toprule
Method & SHD & Precision & Recall & $F_1$ & CPT TV \\
\midrule
Hill climbing, BIC     & \textbf{1.13} & 0.987 & 0.920 & \textbf{0.950} & 0.0110 \\
Exhaustive, BIC        & \textbf{1.13} & 0.987 & 0.920 & \textbf{0.950} & 0.0110 \\
Granger, BIC           & 2.40 & 0.942 & 0.895 & 0.913 & 0.0172 \\
DYNOTEARS              & 2.80 & \textbf{1.000} & 0.778 & 0.866 & 0.0284 \\
n$^2$DBN, selected     & 3.27 & 0.869 & 0.762 & 0.806 & 0.0239 \\
Hill climbing, BDeu    & 4.60 & 0.682 & \textbf{0.967} & 0.778 & 0.0112 \\
Chow-Liu tree          & 7.20 & 0.776 & 0.519 & 0.610 & 0.0582 \\
Empty graph            & 11.60 & 0.000 & 0.000 & 0.000 & 0.1054 \\
\bottomrule
\end{tabular}
\end{table}

The ordering is not the same as on the synthetic grid, and the reason is a
property of the two benchmarks rather than a contradiction. A tiled static
network has no persistent latent state, so the confound that
\cref{sec:results:comparison} isolates is absent and the classical criteria
have nothing to be confounded by. Hill climbing under BIC leads at $F_1 =
0.950$. The local score, with its free quantities selected rather than
declared, reaches $0.806$, up from $0.532$ under the declared setting of the
first version of this experiment.

The per-network breakdown locates what is left. On asia the local score matches
hill climbing exactly, $F_1 = 0.933$ for both, and on earthquake both recover
the network exactly, $F_1 = 1.000$ at zero structural distance. On child, the
largest network here at $20$ nodes, it reaches $0.870$ against $0.933$ at a
higher precision, $0.958$ against $0.933$. On sachs it reaches $0.812$ against
$0.980$. The one clear loss is cancer, $0.413$ against $0.905$, a five-node
network with four edges, where the finest-scale allocation has very little to
condition on. Only three true edges separate the two learners across both
networks; \cref{sec:lim_recall} names them and reports what is known about why,
having tested and withdrawn the block-structure explanation we first proposed.

Conditional probability recovery follows structure recovery and separates the
methods less. The local score reaches a mean total variation distance of
$0.0239$ per conditioning row against $0.0110$ for hill climbing and $0.1054$
for the empty graph, so the tables it implies are far closer to the released
ones than reporting no dependence, at about twice the distance of the best
structural entry. Table~\ref{tab:results:cpt} reports the floor the sample
size allows, estimating each table on the true graph.

\begin{table}[t]
\centering
\caption{Conditional probability table recovery against the released tables
given the true structure, which is the floor the sample size allows rather than
a method score. Mean total variation distance per conditioning row. A node fit
is one node on one replicate. Source \texttt{results/exp8/exp8\_cpt.csv}.}
\label{tab:results:cpt}
\begin{tabular}{lrrrr}
\toprule
Network & Node fits & Mean & Median & Max \\
\midrule
asia       & 24 & 0.0198 & 0.0078 & 0.1317 \\
cancer     & 15 & 0.0122 & 0.0057 & 0.0671 \\
child      & 60 & 0.0194 & 0.0176 & 0.0457 \\
earthquake & 15 & 0.0315 & 0.0278 & 0.1160 \\
sachs      & 33 & 0.0240 & 0.0245 & 0.0465 \\
\midrule
All        & 147 & 0.0210 & 0.0171 & 0.1317 \\
\bottomrule
\end{tabular}
\end{table}

Given the true structure, $94.6$ percent of the $147$ node fits fall below a
total variation distance of $0.05$ and $98.6$ percent below $0.10$, so the
sampled data at this length identifies the released tables well and the
differences in the last column of Table~\ref{tab:results:external} are
attributable to structure rather than to estimation noise.

\subsection{Selecting the free quantities of the model}
\label{sec:results:selection}

A latent-variable model carries free quantities that a latent-free learner does
not. The classical learners in Table~\ref{tab:results:external} each choose
their own penalty through their own score, while the number of regimes $M_k$,
the factor dimension $r_k$ and the edge penalty $\lambda_k$ were held at
declared values. A table produced that way compares declarations, not criteria.
All three are therefore selected on held-out log evidence, computed on the
final quarter of the training series and never on the evaluation data, with the
graph budget and the block lengths of the hierarchy left as declarations
because no held-out score supplies a question about which lags and which
timescales are of interest.

Table~\ref{tab:results:frontier} sweeps $\lambda_k$ on both benchmarks and
reports $F_1$ beside the held-out negative log evidence.

\begin{table}[t]
\centering
\caption{Edge-penalty frontier on both benchmarks, with the held-out negative
log evidence of the same fit. External columns are four tiled networks (asia,
cancer, earthquake, sachs) at $T = 4000$ and block length $63$; synthetic
columns are all $22$ scenarios of the grid at the true hierarchy. Every cell
is read back from \texttt{results/exp9/exp9\_measurements.csv}. The minimum of
the evidence column and the maximum of every $F_1$ column fall on the same row
on the external networks and on different rows on the grid. The grid carries
true edges at the middle and the coarsest scale, so both are reported, together
with the mean over all three scales.}
\label{tab:results:frontier}
\small
\begin{tabular}{lrrrrrr}
\toprule
& \multicolumn{2}{c}{External} & \multicolumn{4}{c}{Synthetic} \\
\cmidrule(lr){2-3}\cmidrule(lr){4-7}
$\lambda_k$ & $F_1$ & held-out NLL & $F_{1,k=2}$ & $F_{1,k=3}$ & $F_{1,\mathrm{mean}}$ & held-out NLL \\
\midrule
$2.0$  & 0.666 & $-68061.4$ & \textbf{0.685} & 0.833 & \textbf{0.597} & $-37890.7$ \\
$1.0$  & 0.725 & $-68065.5$ & 0.669 & 0.940 & 0.582 & $-37923.6$ \\
$0.5$  & 0.812 & $-68067.6$ & 0.602 & \textbf{0.955} & 0.519 & $\mathbf{-37958.6}$ \\
$0.25$ & \textbf{0.869} & $\mathbf{-68070.9}$ & 0.549 & 0.857 & 0.468 & $-37923.8$ \\
$0.1$  & 0.791 & $-68068.9$ & 0.589 & 0.731 & 0.440 & $-37898.8$ \\
\bottomrule
\end{tabular}
\end{table}

Held-out evidence selects the maximizing penalty on the external networks and
does not on the grid. Externally it picks $\lambda_k = 0.25$, which is also
where $F_1$ peaks, and the evidence turns back up on exactly the row where
$F_1$ turns back down. On the grid it picks $0.5$ while $F_{1,k=2}$ peaks at
$2.0$, and the miss is not a tie: $F_{1,k=2}$ is $0.602$ at the selected
penalty against $0.685$ at the maximizing one. It is worse than that on the
columns $F_{1,k=2}$ hides. The strictest setting, $\lambda_k = 2$, reaches a
scale-$2$ precision of $0.864$ and a total structural distance of $6.77$ against
$0.520$ and $12.82$ at the selected $0.5$. Held-out evidence therefore prefers a
graph that is clearly worse at the middle scale.

The frontier was previously scored at the middle scale alone, which raised the
possibility that the disagreement was an artifact of reading a two-scale truth
through a one-scale metric. Rescoring it across both truth-bearing scales
removes that possibility rather than narrowing it to one scale. Held-out
evidence picks $0.5$; $F_{1,k=3}$ peaks at $0.5$ as well, while $F_{1,k=2}$ and
$F_{1,\mathrm{mean}}$ peak at $2.0$. So the selector agrees with the coarsest
scale and disagrees with the middle one, which is a sharper statement than the
blanket disagreement an earlier version of this grid produced and points at the
same conclusion.

That conclusion is the one the wider reading forces. The two truth-bearing
scales do not agree with each other: the middle scale is best at the strictest
penalty and the coarsest at an intermediate one, so no single $\lambda$
maximizes both and a scalar penalty cannot be optimal at every scale on this
grid, whatever selects it. And $F_{1,\mathrm{mean}}$ is not the neutral arbiter
it looks like, since the finest scale is scored for silence and its contribution
recovers as the penalty rises, so the mean rewards a strict penalty partly for
suppressing false positives rather than for finding true edges. The external
half of the table is unaffected throughout, since those networks are tiled flat
and scored at one scale by construction.

The reason is visible elsewhere in these results and is not special pleading.
Held-out likelihood does not order the methods the way structural distance does
on this grid, at a Spearman correlation of $-0.333$ between the two orderings.
A spurious lagged self edge is a genuinely good predictor, because the
autocorrelation it absorbs is real, so a criterion that rewards prediction
prefers denser graphs than the truth wherever a persistent latent state drives
the observations. The tiled networks have no such state, which is why the same
rule works there. The honest statement is therefore conditional: selecting the
penalty on held-out evidence is a principled rule that uses no ground truth,
and it recovers the maximizing setting when prediction and structure agree,
which is the case without a latent confound and not with one. That is a
limitation of held-out evidence as a selector rather than of the score being
selected, and it means an analyst who wants precision on data with a latent
regime should declare a strict penalty and disclose it rather than select one.

\paragraph{The imposed block length is not incidental.}
A tiled network is flat, so the second scale is the analyst's construction. The
finest-scale term of the score is a within-block allocation, so the evidence it
carries about a lagged effect grows with the block length while the number of
blocks falls. Measured at $\lambda_k = 0.5$ on the four networks, $F_1$ rises
from $0.511$ at block length $2$ through $0.601$ and $0.763$ at lengths $5$ and
$10$ to $0.812$ at length $20$, and is unchanged at $40$ and $80$. The block length is declared as the square root of the series length,
which lands inside that plateau and is fixed before any graph is scored.

\paragraph{Subsample stability is measured and not adopted.}
Selection over five contiguous subsamples in the sense of
\citet{meinshausen2010stability} raises external precision to $1.000$ at every
threshold tried, $0.6$, $0.8$ and $1.0$, but costs recall, giving $F_1$ of
$0.831$, $0.812$ and $0.802$ against $0.869$ for the single fit. Held-out
evidence agrees with $F_1$ here too, preferring the single fit at $-68070.7$
against $-68066.7$ for every threshold, so the extra layer is reported and left
out.

\subsection{The search is not what limits recovery}
\label{sec:results:search}

Greedy forward-backward accepts one edge at a time, so a pair of edges that
pays only jointly is unreachable from the incumbent set. Because the score
decomposes per target and every candidate is lagged
(Proposition~\ref{prop:learn_decomposition}), the per-target optimum can be certified
by enumeration whenever the in-degree bound keeps the subset count small, which
it does on both benchmarks: $232$ parent sets for an eleven-node network at
in-degree three, $2626$ for the grid. An exhaustive search and a genetic search
over parent sets were both implemented and run in place of the greedy one.

All three return the same parent sets on every cell of both benchmarks, so
every structural column agrees exactly: $F_1$ of $0.725$ and $0.812$ externally
at $\lambda_k$ of $1.0$ and $0.5$, and $F_{1,k=2}$ of $0.647$ and $0.546$ on
the grid at $2.0$ and $1.0$. The grid figures here run over the four-scenario
subset this comparison uses, which is why they sit below the corresponding rows
of \cref{tab:results:frontier}, an average over every scenario. The held-out losses differ in the fourth
significant figure, since the evolutionary search draws its own randomness and
the fitted coefficients are not bit-identical, but the graphs are. The greedy answer is therefore already
the certified per-target optimum on these data, and no search can improve it.
What limits recovery is the score and the protocol around it, not the
optimizer, which is why the preceding subsection is where the gain is.

\subsection{Aggregation consistency and the recorded objective}
\label{sec:results:invariants}

Two properties are checked on every fit rather than sampled. The first is that
the learned model reconciles across scales, so that for every coarse epoch $c$
at scale $k+1$ and every label $\ell$,
\begin{equation}
X_c^{(k+1,\ell)} \;=\; \sum_{t \in \child_k(c)} X_t^{(k,\ell)},
\qquad
Y_c^{(k+1,\ell)} \;=\; \sum_{t \in \child_k(c)} Y_t^{(k,\ell)} .
\label{eq:results:reconcile}
\end{equation}
The reported reconciliation error is the absolute violation of
Equation~\ref{eq:results:reconcile} summed over coarse epochs, labels and
adjacent scale pairs, so any discrepancy anywhere makes it positive. Across all
$1176$ scored fits of the five grid experiments, and separately on the observed
log below, it is exactly $0$ for both counts and durations, on every row, for
every method including the ablations and the eight external learners. The
equality holds by construction of the allocation and is verified rather than
assumed.

The second is that the recorded objective never decreases. Over the same five
experiments, $1056$ fits recorded a history of at least two outer sweeps,
totaling $4534$ successive steps. Not one of them is a decrease, and the
largest recorded step-to-step change in the wrong direction is exactly $0$. The
rejected-sweep counter, which records a proposal the objective refused, averages
between $0.94$ and $1.72$ per fit across the five experiments and is positive on
$826$ of the $1176$ fits, so the guard is active rather than vacuous.

\paragraph{An observed log.}
The model was also fitted once to an observed event log with no ground truth
(\texttt{results/exp6/exp6\_results.json}), $1601$ events over three labels, one
minute per finest slot, branching $10$ then $6$, so $T_1 = 1440$, $T_2 = 144$
and $T_3 = 24$. Both reconciliation errors are exactly $0$, the recorded
objective rises over all five sweeps, from $7002.16$ to $25044.83$, and one
sweep was rejected. The structure search selected no edge at any scale. The
subsample stability check, which refits on four resamples and records how often
each edge is reselected, returns one candidate at scale $3$, the lag-one edge
from the second label to the third, selected on all four resamples but not on
the full fit. That disagreement between the full fit and the resamples is
reported as it stands; on a log of this length it is a signal to collect more
data, not an edge.

\subsection{An observed log with no ground truth}
\label{sec:results:arm1}

The three preceding groups of tables score recovery against a truth the
generator wrote or a manifest a benchmark declared. This subsection reports the
one arm with neither. The log is LifeSnaps
(\texttt{data/observed\_lifesnaps\_event\_log.csv}), a wearable and
self-report study whose participants enter on their own dates. The panel is
each participant's first eight complete calendar weeks measured from that
participant's own first Monday, which all $51$ of the kept cohort clear, with a
mean observed share of $0.985$ and a minimum of $0.88$. Six behavior labels are
targets, nine columns are exogenous, and the four scales carry $80\,640$,
$1344$, $56$ and $8$ epochs. Three survey columns collapse to one node, since
\texttt{notification\_read} equals the sum of the two survey columns on
$0.971$ of participant-days and no within-group edge is expressible after the
collapse. No claim below is causal. What is claimed is that a graph priced on
this log reaches every scale, that the log drawn back out of the fitted model
reproduces the panel's marginal rates and its coarse-scale non-stationarity,
and that the empty graph does not.

Two stored runs carry the arm. The fit run
(\texttt{results/exp15/exp15\_results.json}) searches one graph over the panel
and grades it on a falsification ladder. The round-trip run
(\texttt{results/exp16/exp16\_results.json}) refits the same search, draws five
synthetic logs per leg from the fitted model, writes them in the raw event
schema, reads them back through the same panel reader and compares them with
the panel on an enumerated grid of non-stationarity forms. The two runs accept
the same $23$ edges, edge for edge, so the logs are drawn from the graph the
ladder grades.

\paragraph{Both directions of the pass are used.}
\Cref{prop:model_block_invariance} splits every per-scale parameter block into
the component a scale identifies and the component it cannot, and the two
directions of the learning pass carry one component each. A search that only
ever bought at one scale would leave that split untested on real data.
Table~\ref{tab:results:arm1_scales} reports where the $23$ accepted edges sit.

\begin{table}[t]
\centering
\caption{The accepted graph by scale and by source, on $180$ priced
candidates at an accepted price of $21.31$ nats against a permutation null
whose median is $0.51$ and whose maximum is $1604.0$. A self edge reads a
label's own past; an exogenous edge reads a calendar column or the scheduled
notification; a cross-label edge reads a different behavior label. The price
strata are self against everything else, so the first two columns together are
the $12$ the run reports as cross. Source
\texttt{results/exp15/exp15\_results.json}.}
\label{tab:results:arm1_scales}
\small
\begin{tabular}{lrrrr}
\toprule
Scale & Exogenous & Self & Cross-label & Total \\
\midrule
$k=1$, minute & $0$ & $0$ & $0$ & $0$ \\
$k=2$, hour   & $12$ & $4$ & $0$ & $16$ \\
$k=3$, day    & $0$ & $4$ & $0$ & $4$ \\
$k=4$, week   & $0$ & $3$ & $0$ & $3$ \\
\midrule
All           & $12$ & $11$ & $0$ & $23$ \\
\bottomrule
\end{tabular}
\end{table}

The search reaches three of the four scales and the split between them is
clean. Every exogenous relation is bought at the hour scale and every edge
bought at the day and week scales reads a label's own past. That is the pattern
the block decomposition predicts. A calendar node is a deterministic function of
the epoch, so by \cref{cor:model_calendar_home} it is identified only where it
varies inside a parent block, and the hour-of-day harmonics have their home
scale where a day is the parent. The scheduled notification is exogenous for a
different reason, that the study delivers it on a schedule the participant does
not set so no behavior precedes it, and it enters at the hour scale because that
is where it varies inside a parent block. What the day and week scales carry
is persistence, which is the component the fine-to-coarse direction supplies and
which a fit at the hour scale alone cannot express. So the two directions of the
pass each buy something on this log, and neither buys what the other does.

The third column of Table~\ref{tab:results:arm1_scales} is the result the
section is least able to soften. \textbf{No edge between two distinct behavior
labels is bought at any scale.} The candidate space carries them and the price
does not exclude them, and none survives the search. The declared anchor set
cannot adjudicate this,
because every positive it declares is the scheduled notification into the survey
node, which is exogenous into a target and not a relation between two behaviors.
Two readings are available and this panel does not separate them. Eight weeks of
$51$ participants may not price a between-label relation above a null whose
maximum is $1604$ nats when the same variance is available from the clock and
from the label's own past, which is where
\cref{sec:results:comparison} shows a latent-free criterion spending its budget.
Or these six labels are conditionally independent given the clock, the
participant level and their own histories, which is a substantive claim about
this cohort and one this arm has no instrument for. The section reports the
count and declines to choose.

The finest scale is expressible and priced and this log does not buy it. A
lag-one self column at the minute scale is identically zero on every scorable
row for four of the six labels, because the risk set excludes exactly the slots
where it would be one, and it is nonzero on $10$ rows of $3\,492\,956$ for a
fifth. Those columns are refused at pool construction on their measured
information, so the price is read on the pool the search actually walks.
A read of one label by another inside the same epoch is admitted at the finest
scale, where one epoch is one slot, and none is bought at $51$ participants,
which repeats what a smaller panel showed. Every lag-zero edge in the tables
above has a calendar source, which is a deterministic function of the epoch and
is never drawn. A same-epoch association is symmetric in its two orientations, so the
target order alone decides which way such an edge is written, and any claim
about direction has to come from outside this fit.

\paragraph{The falsification ladder.}
With no truth to score against, the graph is graded by what it declines to
select. Table~\ref{tab:results:arm1_rungs} reports the rungs, each with its own
denominator printed, since a rate on six is not a performance number.

\begin{table}[t]
\centering
\caption{The falsification rungs of the fit run. R1 and R2 are gates the graph
must not trip; R3 and R9 are read as counts. The declared anchor set is frozen
and hashed before the first scoring run. Source
\texttt{results/exp15/exp15\_results.json}.}
\label{tab:results:arm1_rungs}
\small
\begin{tabular}{llrl}
\toprule
Rung & What it counts & Value & Reading \\
\midrule
R1 & Declared protocol negatives selected & $0/6$ & gate held \\
R2 & Sham-label candidates selected & $1/362$ & $0.003$ false positive rate \\
R3 & Declared positives recovered & $2/6$ & $2$ of the $3$ at the hour scale \\
R9 & Edges reselected on $\ge 0.75$ of $16$ half-panels & $8/23$ & bound of $0.66$ on false ones \\
R7 & Invariant parent subsets for \texttt{gym} at the hour scale & $0/1$ & model rejected \\
\bottomrule
\end{tabular}
\end{table}

R1 holds. Of the $78$ declared protocol negatives, $72$ are structural, because
the scheduled label is an exogenous source and the axis cannot express an edge
into it, so the rung stands on the remaining six and none is selected. R2 is
the informative gate. A sham label is connected to nothing by construction and
one of its $362$ candidates is bought, a false positive rate of $0.003$ inside
the behavioral block. R3 recovers two of the six declared positives. Three of the
six are declared at the minute scale and three at the hour scale, at lags one to
three, and the two recovered are the hour-scale lags one and two; the run also
buys the hour-scale lag zero, which the protocol does not declare. On six the
binomial standard error at one half is $0.204$, so R3 is a gate with its
denominator printed and never a performance number.

R9 is what separates the receipt table from the accepted set. Refitting on
$16$ complementary half-panels, eight of the $23$ edges are reselected on at
least three quarters of them, where the complementary-pairs bound puts at most
$0.66$ low-probability candidates above that threshold. The three
edges with the largest standardized effects in the fit sit among the five with
the lowest reselection rates. Two lag-two self edges into \texttt{nap} at the
hour scale carry standardized effects of $19.2$ and $-145.0$ and are reselected
on one of $16$ half-panels each; their columns clear the support floor and do
not carry enough information to determine what they are given. A coefficient of that size
is a report that its column is nearly empty, and the stability rung is what
removes it. R7 rejects the model for \texttt{gym}. That target carries one
accepted parent at the hour scale, and no subset of it leaves a residual
distribution invariant across participant waves, so the intersection the rung
reports is empty. On one column the rung is a weak test and it is reported with
its width printed.

\paragraph{A calendar node above its home scale buys nothing.}
\Cref{cor:model_calendar_home} is a statement about the kernel and it has a
direct empirical test. Table~\ref{tab:results:arm1_calendar} runs the search
three times on the same panel, with the calendar nodes withheld, offered at
their home scale, and offered at every scale.

\begin{table}[t]
\centering
\caption{The calendar-node ablation, three pre-declared arms on one panel.
The objective is the penalized log-likelihood the search maximizes, so a
larger value is better. Source \texttt{results/exp15/exp15\_results.json}.}
\label{tab:results:arm1_calendar}
\small
\begin{tabular}{lrrrr}
\toprule
Calendar nodes & Candidates & Edges & Objective & vs.\ home \\
\midrule
Withheld          & $172$ & $21$ & $-62\,217.47$ & $-2295.64$ \\
At their home scale & $180$ & $23$ & $\mathbf{-59\,921.83}$ & $0.00$ \\
At every scale    & $204$ & $25$ & $-59\,955.73$ & $-33.90$ \\
\bottomrule
\end{tabular}
\end{table}

Offering the calendar nodes at their home scale is worth $2295.64$ nats over
withholding them. Offering them at every scale costs $33.90$ nats against
offering them at their home scale, while buying two more edges. Both halves of
the corollary hold on this log: the nodes are worth their price where they are
identified, and the price paid above their home scale is wasted. The corollary
turns what is usually a modeling convention into a consequence of the kernel,
and this is the measurement that says the convention is the right one.

\paragraph{The drawn log reproduces the panel's marginal rates.}
The round trip refits the search, draws from the fitted model over the same
roster, mask and window, writes raw events and reads them back. Reconciliation
between counts and durations at adjacent scales is exactly $0$ on every leg and
every replicate, as it is on the synthetic grid, and every drawn log reads back
onto the same $51$ participants and the same observed slots it was drawn over.
Table~\ref{tab:results:arm1_volume} reports the quantity the grid does not
read. Every marginal form of the grid but one divides by the series mean, and
the one exception reads the share of empty epochs, so a leg drawing a label at
half its real rate can still pass every cell it is tested at.

\begin{table}[t]
\centering
\caption{Episode counts on the real panel against the mean over five drawn
replicates, restricted to the same observation mask. The drawn column is the
fitted leg; the last three columns are the drawn-to-real ratio on the fitted
leg, on the leg with the calendar nodes withheld, and on the empty graph. The
summary row is the mean absolute log ratio over the four labels the panel
supports well, which is symmetric in the direction of a miss. Source
\texttt{results/exp16/exp16\_results.json}.}
\label{tab:results:arm1_volume}
\small
\begin{tabular}{lrrrrr}
\toprule
 & & & \multicolumn{3}{c}{Drawn-to-real ratio} \\
\cmidrule(lr){4-6}
Label & Real & Drawn & Fitted & No calendar & Empty \\
\midrule
\texttt{gym} & $506$ & $498.6$ & $0.99$ & $1.00$ & $1.14$ \\
\texttt{nap} & $214$ & $218.4$ & $1.02$ & $1.12$ & $1.24$ \\
\texttt{sleep} & $1841$ & $1919.2$ & $1.04$ & $1.08$ & $1.12$ \\
\texttt{walk} & $1243$ & $1321.0$ & $1.06$ & $0.95$ & $1.10$ \\
\texttt{mindfulness} & $68$ & $132.2$ & $1.94$ & $1.80$ & $1.18$ \\
\texttt{survey\_response} & $10\,162$ & $5418.4$ & $0.53$ & $0.57$ & $0.56$ \\
\midrule
Mean $|\log$ ratio$|$, top four & & & $\mathbf{0.034}$ & $0.061$ & $0.137$ \\
\bottomrule
\end{tabular}
\end{table}

Four of the six labels land within $6$ percent of their real volume. The two
that do not are the two the panel supports least well, for different reasons.
\texttt{mindfulness} occurs $68$ times over $2688$ observed participant-days,
which is one episode per $40$ days, and the draw doubles it. The
\texttt{survey\_response} node carries $72$ percent of the panel's behavioral
episodes and $98.9$ percent of them are zero-duration, so it is a count-only
label whose duration block is uninformative, and the draw reaches $53$ percent
of its rate. Both misses are reported as they stand.

Read as one number, the mean absolute log of the drawn-to-real ratio, which
charges a miss the same in either direction, the volumes are where the graph
separates from the empty one. Over the four well-supported labels the fitted
leg sits at $0.034$, withholding the calendar nodes at $0.061$ and the empty
graph at $0.137$, so the graph is four times closer to the panel than an
intercept and a participant level, and the calendar nodes are half of that
gain. Over all six labels the three legs are $0.239$, $0.232$ and $0.216$, and
the ordering reverses, because the empty graph does not overshoot the label it
has no information about. The two readings are reported together because the
second is the reason for stating which labels the first is read over.

The remaining departure is the share of days left empty. The real panel leaves
$456$ of its $2688$ observed days with no behavioral episode, $17.0$ percent;
the fitted leg leaves $142.8$ per replicate, $5.3$ percent, the no-calendar leg
$133.6$ and the empty graph $89.8$. The ordering is the same one the volumes
give, and no row of the grid reads this quantity.

\paragraph{Non-stationarity is preserved at the coarse scales.}
The comparison is an enumerated grid of $21$ forms over six labels and three
calendar scales, $228$ cells, each an equivalence test at its own declared
tolerance with the participant as the sampling unit. A cell is put to a test
only when the real panel supports one, and $182$ of the $228$ are descriptive
for a declared reason, $127$ of them because the tolerance is inside the
precision the panel can resolve. Table~\ref{tab:results:arm1_grid} reports the
four legs.

\begin{table}[t]
\centering
\caption{The non-stationarity grid, four legs of five replicates each. The
fitted leg carries the searched graph and an hour-scale regime chain; the
no-regime leg is the same search with the chain given one state; the
no-calendar leg is the same search with the clock nodes withheld; the empty
leg carries an intercept and a participant level. Source
\texttt{results/exp16/exp16\_results.json}.}
\label{tab:results:arm1_grid}
\small
\begin{tabular}{lrrrr}
\toprule
Leg & Edges & Tested cells & Failing cells & Failing rows \\
\midrule
Fitted        & $23$ & $46$ & $12$ & $5$ \\
No regime     & $23$ & $46$ & $13$ & $6$ \\
No calendar   & $21$ & $46$ & $12$ & $5$ \\
Empty graph   & $0$  & $46$ & $14$ & $6$ \\
\bottomrule
\end{tabular}
\end{table}

The panel fails the empty graph, which is the control that says it can fail at
all. It also fails the fitted leg, since the verdict is a conjunction over every
tested cell and twelve of the $46$ do not pass.

The row-level reading is the informative one and it needs the positive control
first. Grading one drawn replicate of the fitted leg against the other four,
which nothing separates, five of the $11$ tested rows fail: the
seasonality-amplitude drift, the dispersion, the factor share, the same-epoch
dependence drift and the lagged dependence drift. A row that fails its own
replicate noise has a tolerance narrower than the spread of the draw, so the
fitted leg failing that same row says nothing about the fit. Six rows carry
evidence once those five are removed, and \textbf{the fitted leg passes four of
them}: the normalized trend, the unit-root pair, the share of Haar detail energy
the coarsest level holds, and the windowed maximum mean discrepancy. It fails
the within-day harmonic profile and the drift in persistence.

The same reading bounds what the comparison against the empty graph can carry.
The fitted leg fails two cells fewer of $46$, and the one row that separates
them is the drift in same-epoch dependence, which is one of the five that fail
their own noise. \textbf{On the six rows that carry evidence the two legs are
identical}, passing the same four and failing the same two. So this grid, on
this roster, does not separate a fitted graph from an intercept and a
participant level, and the quantity that does is the volume comparison above,
which the grid does not read.

Table~\ref{tab:results:arm1_cuts} reads the same question as a segmentation.
Both logs are cut independently at their own changepoints under a Poisson cost
and the two cut sets are compared.

\begin{table}[t]
\centering
\caption{Changepoint agreement between the real panel and the first drawn
replicate of the fitted leg, per label and per scale, under a Poisson cost.
Each participant's two series are cut independently and matched, and both
columns are means over the $51$ participants. $F_1$ is scored inside a margin
of three epochs; annotation error is the absolute difference in the number of
cuts. Source \texttt{results/exp16/exp16\_results.json}.}
\label{tab:results:arm1_cuts}
\small
\begin{tabular}{lrrrrrr}
\toprule
 & \multicolumn{2}{c}{Week} & \multicolumn{2}{c}{Day} & \multicolumn{2}{c}{Hour} \\
\cmidrule(lr){2-3}\cmidrule(lr){4-5}\cmidrule(lr){6-7}
Label & $F_1$ & Ann.\ err. & $F_1$ & Ann.\ err. & $F_1$ & Ann.\ err. \\
\midrule
\texttt{gym}             & $0.961$ & $0.06$ & $0.784$ & $0.29$ & $0.765$ & $0.41$ \\
\texttt{mindfulness}     & $0.954$ & $0.06$ & $0.915$ & $0.10$ & $0.890$ & $0.16$ \\
\texttt{nap}             & $0.961$ & $0.04$ & $0.922$ & $0.08$ & $0.902$ & $0.12$ \\
\texttt{sleep}           & $0.784$ & $0.24$ & $0.667$ & $0.41$ & $0.627$ & $0.51$ \\
\texttt{walk}            & $0.706$ & $0.29$ & $0.588$ & $0.49$ & $0.471$ & $0.76$ \\
\texttt{survey\_response} & $0.601$ & $0.53$ & $0.355$ & $1.41$ & $0.054$ & $6.27$ \\
\bottomrule
\end{tabular}
\end{table}

Agreement is highest at the week scale and falls monotonically toward the hour
scale for all six labels. At the week scale five of the six reach $F_1$ above
$0.70$, three of them above $0.95$, and those five miss the panel's cut count by
under a third of a cut. At the hour scale three labels hold above $0.75$ and
\texttt{survey\_response} collapses to $0.054$ at an annotation error of $6.27$
cuts. So the departure grows toward the fine scale for every label, and it is
larger by an order of magnitude for the one whose duration block carries
nothing. That agrees with the grid, whose two evidence-bearing failures are the
within-day harmonic profile and the drift in persistence.

Between-participant dispersion is read as the coefficient of quartile
variation over weekly totals, with a percentile bootstrap over participants.
Three of the six labels are comparable, \texttt{nap} at $0.333$ against
$0.333$, \texttt{walk} at $0.385$ against $0.385$ and \texttt{survey\_response}
at $0.121$ against $0.243$ on overlapping intervals. \texttt{sleep} separates
by $0.021$. \texttt{gym} has overlapping intervals and is called incomparable
because the two panels define the measure for shares of the roster $0.388$
apart. \texttt{mindfulness} is refused: every participant defining a value for
it is pinned at $1.0$ by a zero lower quartile, so the label carries no
dispersion to compare.

\paragraph{What the controls establish and what they do not.}
Four negative controls are declared against the grid. The no-regime control
does its job, breaking every row it is declared against with nothing spurious.
The other three do not. Withholding the calendar nodes leaves the
seasonality-amplitude drift unbroken; shuffling labels across the cross-section
leaves the same-epoch dependence drift unbroken; and a within-participant time
shuffle, which has to break every marginal row, leaves five of them unbroken
and eight untested. A control that fails to break a row it was declared against
says the grid cannot reach that row on this roster, and the honest reading is
that $46$ tested cells of $228$ is the sensitivity this panel supports. That
bound, and not the fitted model, is what the three failing controls measure.
It is also why the row-level reading above is reported net of the positive
control and why the volume table is reported at all.


\subsection{Cost}
\label{sec:results:cost}

Wall-clock cost is reported for the fits as run, one process per fit, with
thread-level parallelism pinned to one so that the numbers are comparable
across methods.

On the synthetic grid a full fit averages $72.1$ seconds with a median of
$60.8$ and a maximum of $203.0$, over a mean of $5.4$ outer sweeps. The empty
graph, which runs the same inference on a fixed graph, averages $14.3$ seconds,
and the oracle-parameter row, which runs inference once and fits nothing,
averages $1.6$. The five grid experiments together account for $96\,886$
seconds of fitting, of which the ten-method comparison of
Section~\ref{sec:results:comparison} is $68\,361$.

The cost is dominated by the structure search rather than by inference; on the
grid the empty-graph row runs the same inference on a fixed graph for $14.3$
seconds against $72.1$ for the full fit. On the external benchmark the local
score costs $174.7$ seconds per fit on average and $461.7$ at most, against
$6.6$ for exhaustive search, $0.3$ for either hill climbing variant, $0.02$
for Granger and $0.003$ for DYNOTEARS. Most of that difference is the selection of
\cref{sec:results:selection}, which fits ten candidate configurations before
the one it keeps, so a single fit at a declared setting costs about a tenth of
the reported figure. Even so the local score is between twenty six times and
five orders of magnitude more expensive than the classical entries, exhaustive
search being the nearest of them and DYNOTEARS the furthest, and its position
in Table~\ref{tab:results:external} does not depend on that. Cost is not the
reason to prefer any row in these tables, and it is reported so that the
comparison is not read as free.
